\documentclass[10pt]{article}
\usepackage[T1]{fontenc}
\usepackage[utf8]{inputenc}
\usepackage[french]{babel}
\usepackage{amsmath,amssymb,mathrsfs,mathpazo}
\usepackage{geometry}
\geometry{a4paper,margin=2cm}
\usepackage{graphicx,array,tabularx,tabularray,colortbl,multirow,multicol,booktabs,enumitem,pifont,ulem,fancybox,diagbox,cancel}
\usepackage{pstricks,pstricks-add,pst-plot,pst-tree,pst-node,pst-text,pst-func,pst-math,pst-bspline,pst-3dplot,pst-eucl,pst-solides3d}
\let\pscalculate\undefined
\usepackage{pst-bezier}
\usepackage{tikz,pgfplots,tkz-tab}
\usetikzlibrary{arrows,arrows.meta,patterns,calc,positioning,trees,decorations.markings,shapes.geometric,shapes.misc,shapes.arrows,intersections,decorations.pathmorphing,decorations.pathreplacing,backgrounds,shadows,babel}
\pgfplotsset{compat=1.18}
\usepgfplotslibrary{fillbetween}
\usepackage[np]{numprint}
\usepackage{listings,makecell,pas-tableur,twoopt}
\let\barre\undefined
\usepackage[pstricks]{tablvar}
\definecolor{darkgreen}{RGB}{0,100,0}
\providecommand{\red}{\color{red}}
\providecommand{\blue}{\color{blue}}
\providecommand{\green}{\color{green}}
\providecommand{\black}{\color{black}}
\providecommand{\white}{\color{white}}
\providecommand{\gray}{\color{gray}}
\providecommand{\cyan}{\color{cyan}}
\providecommand{\magenta}{\color{magenta}}
\providecommand{\yellow}{\color{yellow}}
\providecommand{\violet}{\color{violet}}
\providecommand{\orange}{\color{orange}}
\providecommand{\euro}{\texteuro}
\providecommand{\EUR}{\texteuro}
\providecommand{\ieme}{\textsuperscript{e}}
\providecommand{\degres}{\ensuremath{^\circ}}
\providecommand{\degree}{\ensuremath{^\circ}}
\providecommand{\vv}[1]{\overrightarrow{#1}}
\providecommand{\Coord}[1]{#1}
\providecommand{\decosix}{}
\providecommand{\starredbullet}{\ensuremath{\bullet}}
\pagestyle{empty}
\setlength{\parindent}{0pt}
\usepackage{adjustbox}\DeclareUnicodeCharacter{2236}{:}\DeclareUnicodeCharacter{2190}{\ensuremath{\leftarrow}}
\providecommand{\R}{\mathbb{R}}
\providecommand{\N}{\mathbb{N}}
\providecommand{\Z}{\mathbb{Z}}
\providecommand{\Q}{\mathbb{Q}}
\providecommand{\D}{\mathbb{D}}
\providecommand{\C}{\mathbb{C}}
\providecommand{\e}{{\mathrm{e}}}
\providecommand{\vect}[1]{\overrightarrow{#1}}
\providecommand{\vectt}[1]{\overrightarrow{\mathrm{#1}}}
\providecommand{\vvt}[1]{\overrightarrow{\mathrm{#1}}}
\providecommand{\np}[2][]{{#2~#1}}
\providecommand{\vv}[1]{{\overrightarrow{#1}}}
\renewcommand{\d}{{\mathrm{d}}}
\providecommand{\degres}{{^\circ}}
\providecommand{\degre}{{^\circ}}
\providecommand{\blue}{\color{blue}}
\providecommand{\red}{\color{red}}
\providecommand{\enskip}{\quad}
\providecommand{\llbracket}{\lbrack\!\lbrack}
\providecommand{\rrbracket}{\rbrack\!\rbrack}
\providecommand{\euro}{{\text{€}}}
\providecommand{\no}{{\text{n°}}}
\providecommand{\interval}[2]{{[#1;#2]}}
\providecommand{\intervFO}[2]{{[#1;#2[}}
\providecommand{\dotfill}{\ldots}

\providecommand{\ds}{\displaystyle}
\providecommand{\textding}[1]{#1}
\providecommand{\Oij}{\(\mathrm{O};\vec{\imath},\vec{\jmath}\)}
\providecommand{\Oijk}{\(\mathrm{O};\vec{\imath},\vec{\jmath},\vec{k}\)}
\providecommand{\Ouv}{\(\mathrm{O};\vec{u},\vec{v}\)}
\providecommand{\intervFF}[2]{\ensuremath{[#1;#2]}}
\renewcommand{\euro}{\ifmmode\text{\texteuro}\else\texteuro\fi}
\DeclareUnicodeCharacter{00B0}{\ifmmode^{\circ}\else\textdegree\fi}
\DeclareUnicodeCharacter{20AC}{\ifmmode\text{\texteuro}\else\texteuro\fi}
% Énoncés : Éducation nationale. Transcription et corrigés : APMEP, Denis Vergès et contributeurs. Découpage et classement : CoopMaths, Sébastien Lozano, Éric Elter et contributeurs. Adaptation de la présentation : Maths Explorer.
% CC BY-SA : https://coopmaths.fr/www/about/
% https://coopmaths.fr/alea/static/bac/2026/tex/bac_2026_06_sujet2_etrangers_4.tex
\begin{document}
\noindent\textit{APMEP / CoopMaths. CC BY-SA. Présentation adaptée par Maths Explorer.}\par\medskip
\subsection*{Partie A}
Sur le graphique ci-dessous, on a tracé trois courbes $\mathcal{C}_{1}$\,, $\mathcal{C}_{2}$ et $\mathcal{C}_{3}$\,.

Les courbes correspondent aux représentations graphiques de trois fonctions définies sur $\R$ : une fonction $f$, sa dérivée $f'$ et sa dérivée seconde $f''$.

\begin{center}
	\begin{tikzpicture}
		\begin{axis}[x={17mm},y={70mm},
			xmin= -0.4,  xmax= 6.1, ymin = -.6, ymax=1.15,
			xtick distance=1, ytick distance=0.5,
			tick label style={font=\footnotesize}, grid=major, axis lines =center]
			\addplot [line width=1.2pt,color=red,smooth,samples=100,domain= 0:6.1,dotted]{exp(-x)*(x-1)*(x-2)};
			\addplot [line width=1.2pt,color=teal,smooth,samples=100,domain= 0:6.1,dashed]{exp(-x)*(-x*x+5*x-5)};
			\addplot [line width=1.2pt,color=blue,smooth,samples=100,domain= 0:6.1]{exp(-x)*(x*x-7*x+10)};
			\node[below left=2pt] at (0,0) {{\footnotesize 0}};
			\node[red] at(0.5,0.9) {$ \mathcal{C}_1 $};
			\node[blue] at(1.4,0.95) {$ \mathcal{C}_2 $};
			\node[teal] at(0.75,-0.45) {$ \mathcal{C}_3 $};
		\end{axis}
	\end{tikzpicture}
\end{center}

Associer chacune des fonctions $f$, $f'$ et $f''$ à sa courbe représentative. \emph{Aucune justification n'est attendue.}


\subsection*{Partie B}
On considère l'équation différentielle $(E)$ définie par $y'+y=(2 x-3) \e^{-x}$ où $y$ est une fonction de la variable réelle $x$.

\begin{enumerate}
	\item On considère la fonction $g$ définie sur $\R$ par
	\[g(x)=\left(x^{2}-3 x\right) \e^{-x}\]

	Démontrer que la fonction $g$ est une solution particulière de l'équation différentielle $(E)$.

	\item Déterminer l'ensemble des solutions de l'équation différentielle \quad $y'+y=0$.

	\item En déduire l'ensemble des solutions de l'équation différentielle $(E)$.

	\item Déterminer la solution $f$ de l'équation différentielle $(E)$ telle que \quad $f(0)=2$.
\end{enumerate}

\subsection*{Partie C}
On considère la fonction $f$ définie sur $\R$ par
\[f(x)=\e^{-x}\left(x^{2}-3 x+2\right)\]

et on note $\mathcal{C}_{f}$ sa courbe représentative dans un repère orthogonal.

\begin{enumerate}
	\item Étudier le signe de la fonction $f$ sur $\R$.

	\item \begin{enumerate}
		\item Déterminer la limite de la fonction $f$ en $-\infty$.

		\item Déterminer la limite de la fonction $f$ en $+\infty$.
	\end{enumerate}

	\item On note I l'intégrale définie par :
	\[I=\int_{0}^{1} f(x) \d x\]

	\begin{enumerate}
		\item À l'aide de deux intégrations par parties successives, démontrer que \quad $I=1-\dfrac{1}{\e}$.

		\item Interpréter graphiquement ce résultat.
	\end{enumerate}
\end{enumerate}

\subsection*{Partie D}
On considère un réel $a$.

On note $\left(T_{a}\right)$ la tangente à la courbe $\mathcal{C}_{f}$ au point d'abscisse $a$.

\begin{enumerate}
	\item Démontrer que le point d'intersection de la tangente $(T_{a})$ et de l'axe des ordonnées a pour ordonnée
	\[\left(a^{3}-4 a^{2}+2 a+2\right) \e^{-a}\]

	\item Déterminer le nombre de tangentes à la courbe $\mathcal{C}_{f}$ passant par l'origine du repère.

	Le candidat explicitera les étapes de la démarche utilisée.
\end{enumerate}

\end{document}
